% worked_examples.tex — step-by-step runbooks for the checked-in examples/*
% topologies, ordered from the easiest to the most advanced. Companion to,
% not a repeat of, the per-command reference in Text/user_guide.tex and the
% narrative tutorials docs/tutorial_cgsi1.md / docs/tutorial_cawaqs.md: this
% section gives the exact command sequence for each topology and explains
% what changes between them, and points to those other documents for flag
% semantics and full expected-output transcripts.
% -----------------------------------------------------------------------
\section{Worked Examples: From \texttt{CGSil1} to \texttt{cawaqs}}
\label{sec:worked-examples}

\cgspkg{} ships four checked-in \texttt{examples/} topologies. Run through
them in order: each one adds exactly one new idea on top of the last, so
together they form a guided path from the simplest possible \texttt{.cgs}
to a real 19-repository build tree.

\begin{center}
\begin{tabular}{lllp{5.3cm}}
\textbf{Level} & \textbf{Topology} & \textbf{Entry point} & \textbf{New idea introduced} \\
\hline
1 & \texttt{CGSil1} & hand-authored \texttt{.cgs} & minimal spec, name-matching auto-mount \\
2a & \texttt{cawaqsviz} & hand-authored \texttt{.cgs} & explicit overrides on a real, irregular tree \\
2b & \texttt{cawaqsviz} & \texttt{discover} & drafting a \texttt{.cgs} from a checkout, no authoring at all \\
2c & \texttt{cawaqsviz} & \texttt{import-submodules} & migrating an existing git-submodules project in \\
3 & \texttt{cawaqs} & hand-authored \texttt{.cgs} & the same minimal-field style as Level~1, at 19-repo scale \\
\end{tabular}
\end{center}

Levels~2a--2c are not three different projects: they are three different
\emph{starting points} for the same real project, \texttt{cawaqsviz}, chosen
so this section can demonstrate every route \cgspkg{} offers for arriving
at a working \texttt{.cgs} --- write it by hand, derive it from a checkout,
or migrate it from an existing submodules setup --- without inventing a
fourth toy topology to do it.

% -----------------------------------------------------------------------
\subsection{Level 1 --- \texttt{CGSil1}: the minimal synthetic topology}

\texttt{CGSil1} is the easiest way into \cgspkg{}: a 4-repository,
mixed-provider tree with every field left at its default. Its full
\texttt{.cgs} and the rationale for each line are in \S\ref{chap:user-guide}
``Local Authoring Spec'' above; \texttt{docs/tutorial\_cgsi1.md} carries the
full transcript with expected output for every step. What follows is the
command sequence itself.

\begin{lstlisting}[language=bash]
# 1. Fetch the root repo and cgitsync itself (nested mode: cgitsync is
#    cloned inside the tree it will manage).
$ git clone https://gitlab.com/CGS_test/CGSil1
$ cd CGSil1 && git clone https://github.com/flipoyo/ComplexGitSync
$ cd ComplexGitSync

# 2. Check the spec offline, no network beyond parsing.
$ pixi run cgitsync validate ../CGSil1.cgs

# 3. Preview the topology before touching disk.
$ pixi run cgitsync view-tree ../CGSil1.cgs

# 4. Clone every declared repo and reach READY.
$ pixi run cgitsync initialise ../CGSil1.cgs

# 5. The ordinary git cycle, run tree-wide from here on --- no path
#    argument needed, the .gts snapshot is discovered automatically.
$ pixi run cgitsync pull
$ pixi run cgitsync add
$ pixi run cgitsync commit "my commit message"
$ pixi run cgitsync push

# 6. Minimalist release: stage, commit, pull, push, freeze in one call.
$ pixi run cgitsync freeze-release v1.1.0 "release v1.1.0"

# 7. Return to a frozen release at any later point.
$ pixi run cgitsync launch-release v1.1.0
\end{lstlisting}

The only \texttt{.cgs} idea this level relies on is the
\textbf{name-matching auto-mount}: because \texttt{CGSil1}'s repository
identifier's last segment equals \texttt{project = "CGSil1"}, the root is
inferred at \texttt{relative\_path = "."} with no override written anywhere.
That convenience is exactly what stops being safe once a project's
identifier and display name diverge --- which is where Level~2 picks up.

% -----------------------------------------------------------------------
\subsection{Level 2 --- \texttt{cawaqsviz}: three ways to the same \texttt{.cgs}}

\texttt{cawaqsviz} (\url{https://gitlab.com/cawaqs/gviz/cawaqsviz}) is a
real GitLab project, nested three path segments deep, with two GitHub
children that were, until recently, plain git submodules. It has no
\texttt{.cgs} of its own anywhere upstream. That makes it a good vehicle for
showing all three ways \cgspkg{} can produce a working \texttt{.cgs} for a
project it does not already control the source of, ordered here by how much
the operator has to already know about the tree before running the command.

\subsubsection{2a --- Write it by hand}

The most explicit, and most error-prone, route: read the project's own
\texttt{.gitmodules} and topology, then author \texttt{examples/cawaqsviz.cgs}
directly. Its corrected content and the exact mistakes the first draft made
(wrong nested-subgroup identifier, missing \texttt{relative\_path}
overrides, missing \texttt{nested\_config = "disabled"}) are already spelled
out in \S\ref{chap:user-guide} ``Local Authoring Spec''; this level is about
running it, not re-deriving it:

\begin{lstlisting}[language=bash]
$ pixi run cgitsync validate examples/cawaqsviz.cgs
$ pixi run cgitsync bootstrap examples/cawaqsviz.cgs cawaqsviz-demo
$ pixi run cgitsync view-tree --search-dir <path bootstrap printed>
\end{lstlisting}

\texttt{bootstrap} is used instead of \texttt{initialise} here because, unlike
\texttt{CGSil1} above, this example is not meant to have \cgspkg{} cloned
inside the tree it manages --- see \texttt{bootstrap} in
\S\ref{chap:user-guide} for why the two commands pick different default
locations.

\subsubsection{2b --- Derive it with \texttt{discover} (autodiscover mode)}

The opposite route: start from nothing but a checkout, and let \cgspkg{}
read the filesystem instead of writing the \texttt{.cgs} by hand. This is
the route to prefer whenever a checkout is already available, since it
cannot repeat the hand-authoring mistakes of 2a --- \texttt{relative\_path}
falls out of the walk directly instead of being typed.

\begin{lstlisting}[language=bash]
# A plain clone leaves submodule paths empty; initialise them first so
# discover has something to find at those paths.
$ git clone https://gitlab.com/cawaqs/gviz/cawaqsviz cawaqsviz-scan
$ cd cawaqsviz-scan && git submodule update --init --recursive

# Dry run: report what discover sees, write nothing.
$ pixi run cgitsync discover . --max-depth 3

# Satisfied with the report: draft the .cgs.
$ pixi run cgitsync discover . --write cawaqsviz-discovered.cgs
$ pixi run cgitsync validate cawaqsviz-discovered.cgs
\end{lstlisting}

The draft reconstructs 2a's hand-corrected file almost field-for-field: the
root at \texttt{relative\_path = "."}, both children at their real submodule
paths (\texttt{external/HydrologicalTwinAlphaSeries},
\texttt{docs/CWV\_user\_guide}) rather than their bare repo names, and
\texttt{nested\_config = "disabled"} on both, since neither has a
\texttt{.cgs} of its own --- \texttt{discover} derives that last fact from
the filesystem rather than assuming it, so it cannot omit it the way 2a's
first draft did. See \texttt{discover} in \S\ref{chap:user-guide} for exactly
what is and is not reported (no network calls, unresolvable remotes are
warned about and left out, never guessed at).

\subsubsection{2c --- Migrate it with \texttt{import-submodules}}

The historical route: before either \texttt{.cgs} above existed,
\texttt{cawaqsviz} tracked both children as real git submodules. This mode
starts from that state and converts it, rather than describing a tree that
is assumed to already be plain clones.

\begin{lstlisting}[language=bash]
# Dry run against a real cawaqsviz checkout that still has submodules:
# reports name, path, URL, branch per submodule; changes nothing.
$ pixi run cgitsync import-submodules /path/to/cawaqsviz

# Convert: drop the gitlinks, retire .gitmodules, extend .gitignore,
# and emit the equivalent .cgs entries.
$ pixi run cgitsync import-submodules /path/to/cawaqsviz \
    --apply --output cawaqsviz_submodules.cgs

# Review and commit the resulting working-tree changes as usual.
$ cd /path/to/cawaqsviz && git status
$ git commit -m "chore: retire git submodules in favour of ComplexGitSync"
\end{lstlisting}

The full before/after picture (index, \texttt{.gitmodules},
\texttt{.gitignore}) and the automated test that proves it against fixture
repositories are in \texttt{docs/tutorial\_cgsi1.md} \S6 --- this level is
the condensed command sequence, not a restatement of that walkthrough.
\texttt{import-submodules} and \texttt{discover} answer different questions
about the same tree: \texttt{discover} reports what is \emph{checked out},
\texttt{import-submodules} reads what git \emph{declares} in
\texttt{.gitmodules} and acts on it --- a submodule path that was never
initialised is invisible to 2b but not to 2c.

\subsubsection{A shared caveat across all three}

\cgspkg{} stores no credentials and has no private-repository
authentication story (\S\ref{chap:user-guide} ``\texttt{import-submodules}''
and ``\texttt{discover}''). All three routes above assume every repository
in the tree is reachable anonymously or through credentials the ambient
environment (\texttt{ssh-agent}, an HTTPS credential helper) already holds.

% -----------------------------------------------------------------------
\subsection{Level 3 --- \texttt{cawaqs}: the same minimal style, at scale}

\texttt{cawaqs} is the most advanced example not because it needs more
\texttt{.cgs} features than \texttt{cawaqsviz} --- it needs fewer --- but
because it applies Level~1's minimal, no-override style to a real
19-repository build tree (root + 17 physics libraries across five GitLab
groups + one required build-tooling repo) instead of a synthetic 4-repo one.
Every entry's on-disk layout already matches the default
\texttt{relative\_path = <repo\_name>}, so, unlike \texttt{cawaqsviz}, not a
single override is written anywhere in \texttt{examples/cawaqs.cgs} beyond
the shared branch-fallback convention:

\begin{lstlisting}[language=bash]
$ pixi run cgitsync validate examples/cawaqs.cgs
$ pixi run cgitsync bootstrap examples/cawaqs.cgs cawaqs-smoke-test
$ pixi run cgitsync view-tree --search-dir <path bootstrap printed>

# Move the whole 19-repo tree to a shared branch in one call, falling
# back to main per-repo wherever that branch does not exist.
$ pixi run cgitsync checkout my-feature-branch
\end{lstlisting}

Reaching \texttt{READY} here only replaces the \emph{repository-fetching and
branch-selection} half of \texttt{cawaqs}'s own
\texttt{make\_Cawaqs.sh}/\texttt{make\_Cawaqs\_from\_branches.sh} scripts.
\cgspkg{} does not compile anything: the remaining
\texttt{make -f Makefile} build step is still the user's job, and only
succeeds from this tree with \texttt{LIB\_HYDROSYSTEM\_PATH} left unset or
pointed at the bootstrapped root --- the alternative shared,
decoupled install layout those scripts also support is out of scope for
\cgspkg{}'s containment model. See \texttt{docs/tutorial\_cawaqs.md} for the
full boundary between what \cgspkg{} does and does not own for this project.

\texttt{discover} is not a route into \texttt{cawaqs.cgs} the way it is for
\texttt{cawaqsviz}: the 17 library repositories never coexist inside one
directory tree until \emph{after} a \texttt{.cgs} already lists them, so
there is no single checkout for a filesystem walk to scan. Authoring this
one from documented developer knowledge, as Level~2a does for
\texttt{cawaqsviz}, is the only available route.
